% =====================================================================
%  Entregable 8: Influencia y relevancia en el efecto macro clasico
%  Fragmento \input-able en context/doc/mpemba_cuantico.tex
% =====================================================================
\section{Influencia y relevancia en el efecto macroscópico clásico}
\label{sec:macro}

\subsection{La raíz macroscópica y su problema histórico}
El efecto Mpemba nació como una observación \emph{macroscópica y clásica}: agua
inicialmente más caliente que se congela antes que agua más fría
(Mpemba--Osborne 1969). Durante medio siglo fue un fenómeno disputado, porque a
escala macroscópica intervienen muchos mecanismos difíciles de controlar
(convección, evaporación, sobreenfriamiento, gases disueltos) y porque
\emph{carecía de un criterio teórico limpio}: ¿qué significa exactamente que un
sistema esté ``más lejos del equilibrio''? Distintas métricas de distancia daban
respuestas distintas (la crítica de Burridge--Linden 2016).

\subsection{El mecanismo unificador: geometría espectral de la relajación}
El marco cuántico desarrollado en este trabajo aporta, precisamente, lo que le
faltaba al fenómeno macroscópico: un \textbf{mecanismo espectral riguroso} que es
\emph{el mismo} en el régimen clásico y en el cuántico. En ambos casos la
relajación se descompone en modos,
\[
\text{clásico: } \vb{p}(t)=\pi+\sum_{k\ge2} a_k\, e^{\lambda_k t}\,\vb{v}_k,
\qquad
\text{cuántico: } \rho(t)=\rss+\sum_{k\ge2} a_k\, e^{\lambda_k t}\,\hat r_k,
\]
con $\vb{v}_k$ los autovectores de la matriz de tasas y $\hat r_k$ los
autooperadores del Liouvilliano. El efecto Mpemba ocurre, en \emph{los dos}
mundos, cuando la preparación caliente \emph{solapa poco con el modo lento}
($a_2\approx0$): es un fenómeno de \textbf{geometría no-normal} de los modos de
relajación, no una propiedad del agua. Lo que en el agua era una anomalía
empírica resulta ser un caso particular de una estructura matemática general.

\subsection{Lo que el enfoque cuántico devuelve al problema clásico}
Tres aportaciones del marco cuántico iluminan directamente el efecto macroscópico:
\begin{enumerate}[leftmargin=1.7em,itemsep=2pt]
\item \textbf{Criterio espectral del Mpemba fuerte} ($a_2=0$, Carollo \emph{et
al.}): da una condición \emph{exacta} y constructiva para la aceleración, válida
también para sistemas clásicos de muchos estados (vidrios, granulares).
\item \textbf{Diagnóstico independiente de la métrica} (termomayorización,
Vu--Hayakawa): resuelve la ambigüedad histórica de Burridge--Linden ---certifica
el efecto bajo \emph{todas} las distancias monótonas a la vez--- y es aplicable
tal cual a las distribuciones clásicas.
\item \textbf{Los puntos excepcionales y el recurso metrológico}: conceptos
nacidos en el régimen cuántico que sugieren nuevos protocolos (aceleración por
preparación, termometría fuera del equilibrio) trasladables a plataformas
clásicas controladas.
\end{enumerate}

\subsection{La firma universal: del macro clásico al cuántico}
La figura~\ref{fig:macro} muestra la \emph{misma} firma ---el cruce de las curvas de
relajación, donde una preparación que parte más lejos alcanza el equilibrio
antes--- en tres dominios físicos usando los datos del repositorio: un coloide
clásico en trampa óptica (Kumar 2022), un circuito cuántico de muchos cuerpos
(Turkeshi 2024) y la simulación cuántica propia de este trabajo. La continuidad
conceptual es total: el mismo criterio de solapamiento espectral explica el cruce
en los tres casos.

\begin{figure}[h]
\centering
\includegraphics[width=\textwidth]{../../macro/results/fig_macro_universal.png}
\caption{La firma universal del efecto Mpemba ---el cruce de curvas de
relajación--- del macro clásico al cuántico. (a)~coloide clásico experimental
(Kumar 2022): la distancia al equilibrio de una preparación caliente cruza por
debajo de otra más fría. (b)~circuito cuántico de muchos cuerpos (Turkeshi 2024):
la asimetría de entrelazamiento de $\theta=1{,}2$ parte como la mayor pero se
restaura antes. (c)~simulación cuántica propia (Mpemba fuerte por geometría
espectral). El mecanismo ---solapamiento con los modos lentos--- es común.}
\label{fig:macro}
\end{figure}

\subsection{Relevancia}
La relevancia es doble. \emph{Hacia atrás}: el aparato cuántico
(modos del generador, solapamientos, diagnósticos universales) proporciona al
efecto macroscópico clásico el fundamento teórico riguroso que históricamente le
faltó, convirtiendo una curiosidad sobre el agua en un teorema sobre la
relajación fuera del equilibrio. \emph{Hacia adelante}: muestra que el mismo
fenómeno, y las mismas herramientas computacionales (descomposición espectral,
evolución del operador densidad, trayectorias), abarcan sin costura desde el vaso
de agua hasta el procesador cuántico, lo que justifica tratar el efecto Mpemba
como un \emph{principio general de diseño} de protocolos de relajación acelerada.
